\chapter{Introdução e importância da automação}
\label{cap:introducao}

\begin{objetivos}
Ao final deste capítulo você deve ser capaz de:
\begin{itemize}
  \item explicar o que é automação industrial e por que ela existe;
  \item relacionar os objetivos da automação com os indicadores que ela melhora;
  \item situar o Controlador Lógico Programável (CLP) dentro do sistema de automação;
  \item distinguir, na prática, o que é definido por norma do que é definido por fabricante.
\end{itemize}
\end{objetivos}

\section{O que é automação industrial}

Automação industrial é o uso de equipamentos e programas para operar um processo produtivo com
pouca ou nenhuma intervenção humana direta. Ela não substitui apenas a força física: substitui
sobretudo a \emph{decisão repetitiva} --- ligar, desligar, medir, comparar com um valor desejado e
corrigir.

Essa decisão repetitiva foi, por décadas, responsabilidade de painéis de relés, temporizadores de
came e sequenciadores. A lógica estava \emph{fiada}, literalmente: mudar a sequência de uma máquina
significava refazer o painel elétrico, parar a produção e refazer a documentação. Foi essa rigidez
que criou a necessidade que dá origem ao CLP (Capítulo~\ref{cap:evolucao}).

\index{automação industrial}\index{automação industrial!objetivos}

\section{Por que automatizar}

Automatizar não é um fim em si. As razões aparecem em cinco frentes, e cada uma tem uma forma de
ser medida --- o que importa quando é preciso justificar um investimento
(Figura~\ref{fig:objetivos-automacao}).

\begin{figure}[H]
  \centering
  \includegraphics[width=0.95\linewidth]{01-objetivos-automacao.pdf}
  \caption{Pressões que levam à automação e ganhos que ela entrega.}
  \label{fig:objetivos-automacao}
\end{figure}

\begin{table}[H]
  \centering
  \caption{Frentes de ganho da automação e como observá-las.}
  \label{tab:frentes-automacao}
  \begin{tabular}{p{3.4cm} p{6.2cm} p{4.6cm}}
    \toprule
    Frente & O que muda na prática & Como se observa \\
    \midrule
    Produtividade
      & Menos tempo parado, ciclo mais estável e maior volume no mesmo tempo
      & Disponibilidade, tempo de ciclo, peças por hora \\
    Qualidade
      & Repetibilidade da sequência: a máquina executa igual a cada ciclo
      & Refugo, retrabalho, variação da característica \\
    Segurança
      & Eliminação do homem na zona de risco e intertravamento confiável
      & Incidentes, intervenções em área de risco \\
    Energia e insumo
      & Desligamento automático, ponto de operação otimizado, fim do funcionamento a vazio
      & Consumo específico (kWh por unidade), desperdício de matéria-prima \\
    Informação
      & O processo passa a produzir dados de produção, falha e consumo
      & Existência de histórico, rastreabilidade, indicadores por turno \\
    \bottomrule
  \end{tabular}
\end{table}

\section{Onde o CLP é aplicado}

Praticamente não há ramo industrial sem aplicação de controladores programáveis. Os casos mais
frequentes:

\begin{itemize}
  \item \textbf{Máquinas industriais} --- usinagem, injetoras de plástico, têxteis, calçados;
  \item \textbf{Processos contínuos e de batelada} --- siderurgia, papel e celulose, petroquímica,
        química, alimentos, mineração;
  \item \textbf{Energia} --- controle de demanda e de fator de carga;
  \item \textbf{Supervisão e aquisição de dados} --- fábricas, prédios inteligentes, utilidades;
  \item \textbf{Ensaios e bancadas} --- teste automático de componentes, laboratórios;
  \item \textbf{Sistemas embarcados e móveis} --- máquinas agrícolas e de construção,
        equipamentos de transporte.
\end{itemize}

\section{O CLP no centro do sistema}

O CLP é o equipamento que executa a lógica de controle de forma \emph{cíclica e determinística}:
ele lê as entradas, processa o programa e escreve as saídas, repetidamente, em intervalos de
alguns milissegundos. É essa característica que o diferencia de um computador comum e que dá
previsibilidade ao processo (Capítulo~\ref{cap:funcionamento}).

Ao redor dele ficam os instrumentos de campo, os módulos de entrada e saída, o terminal de
programação e as redes de comunicação que o ligam à supervisão e à gestão
(Capítulo~\ref{cap:estrutura}).

\begin{atencao}
Nada nesta apostila autoriza intervir em painel energizado. Antes de qualquer medição, troca de
módulo ou edição de programa em máquina em produção, aplique o procedimento de bloqueio e
sinalização da sua instituição. CLP não é equipamento de aprendizagem em máquina viva.
\end{atencao}

\section{Como ler esta apostila}

Esta obra foi organizada em cinco partes: fundamentos e projeto, hardware e funcionamento,
programação e normas, integração e sistemas, operação e carreira. Cada capítulo fecha com
exercícios e, quando aplicável, uma caixa \emph{Na prática}.

Uma convenção vale para todo o texto: \textbf{sempre que uma linguagem, instrução ou endereçamento
for apresentado, ficará claro se aquilo é definido por norma ou se é proprietário de um
fabricante}. Essa distinção é o assunto do Capítulo~\ref{cap:iec-versus-legado} e aparece marcada
em cada tópico de programação.

\index{CLP}\index{programação de CLP}\index{norma!IEC 61131-3}

\begin{pratica}
Escolha uma máquina da sua escola ou do seu trabalho e responda por escrito:
(a) quais decisões hoje são feitas por relés, botoeiras ou operador;
(b) quais dessas decisões poderiam ser automatizadas;
(c) qual indicador da Tabela~\ref{tab:frentes-automacao} melhoraria e como você mediria isso hoje,
sem instalar nada novo.
\end{pratica}

\begin{exercicios}
\begin{enumerate}[label=\alph*)]
  \item Explique, com suas palavras, por que ``lógica fiada'' dificulta a mudança de um processo.
  \item Cite dois ganhos da automação que não sejam de produtividade e justifique sua escolha.
  \item Um gestor afirma que automatizar sempre reduz o custo por peça. Aponte uma situação em que
        isso não se verifica e explique o motivo.
  \item Dê um exemplo de decisão que \emph{não} deve ser confiada apenas ao programa do CLP, e diga
        quem ou o que deveria assumi-la.
\end{enumerate}
\end{exercicios}
